% Les styles de tracé de la présentation.
%
% Ils sont écrits ici plutôt qu'empruntés à `report/images.tex` : ce fichier-là porte l'apparat des
% figures du rapport — emplacements réservés, logotypes, compteur de captures — dont une planche
% n'a que faire, et il déclare ses couleurs par des noms que le document principal du rapport
% définit. Une planche qui l'inclurait chargerait un apparat inutile et dépendrait d'un fichier
% qu'elle ne compose pas.
%
% La couleur d'accent est celle de la page de garde du rapport, à l'identique : le support de
% soutenance et le document remis se ressemblent, sans que l'un lise l'autre.

\definecolor{accentgarde}{RGB}{23,54,93}
\definecolor{griscartouche}{gray}{0.42}

\setbeamercolor{structure}{fg=accentgarde}
\setbeamercolor{frametitle}{fg=accentgarde}
\setbeamercolor{title}{fg=accentgarde}
\setbeamerfont{frametitle}{series=\bfseries}
\setbeamertemplate{itemize item}{\color{accentgarde}\textbullet}
\setbeamertemplate{itemize subitem}{\color{griscartouche}--}

\usetikzlibrary{positioning, shapes.geometric, fit, calc, arrows.meta}

\tikzset{
  % — une brique de couche, ou une étape.
  briquePres/.style={
    rectangle, draw=accentgarde, line width=0.6pt, align=center,
    font=\scriptsize, inner sep=2.5pt, minimum height=7mm, text width=24mm,
  },
  briqueLargePres/.style={briquePres, text width=54mm},
  % — un cadre secondaire : trait pointillé.
  briqueOptPres/.style={briquePres, dash pattern=on 1.6pt off 1.6pt},
  % — l'intitulé d'une couche.
  couchePres/.style={font=\scriptsize\scshape, text=griscartouche, align=right},
  % — une association vers une table descriptive.
  assocPres/.style={draw=accentgarde, line width=0.5pt},
  % — une flèche de flux.
  fluxPres/.style={draw=accentgarde, line width=0.6pt, -{Latex[length=1.8mm,width=1.4mm]}},
  % — une accolade de regroupement.
  notePres/.style={font=\scriptsize\itshape, text=griscartouche},
}
